\section{Evaluation Methodology}
\label{sec:evaluation}

\subsection{Experimental Setup and Procedure}

The testbed runs in an AWS Academy environment with one control-plane Node and
six \textit{t3.micro} worker Nodes. Each worker instance provides \(1\,GiB\) of
physical memory. After system and Kubernetes reservations, the recorded
allocatable capacity is approximately \(809\,Mi\) of memory and two CPU cores
per worker.

\begin{table*}[t]
    \centering
    \caption{Recorded Experimental Setup}
    \label{tab:experimental_setup}
    \small
    \renewcommand{\arraystretch}{1.15}
    \begin{tabular}{@{}>{\raggedright\arraybackslash}p{3.1cm}
        >{\raggedright\arraybackslash}p{12.4cm}@{}}
        \toprule
        \textbf{Item} & \textbf{Recorded configuration} \\
        \midrule
        Environment & AWS Academy; one control-plane Node and six
        \textit{t3.micro} workers, each with \(1\,GiB\) physical memory.
        Kubernetes allocatable capacity was two CPU cores and approximately
        \(809\,Mi\) memory per worker. \\
        Runtime components & Kubernetes cluster, Metrics Server through the
        Kubernetes Metrics API, a single-shot Descheduler Job, and
        \textit{k6} running outside the workload Pod through a NodePort. \\
        API workload & \textit{workload-api} Deployment with one replica Pod;
        \(50\,m\) CPU request, \(1000\,m\) CPU limit, \(700\,Mi\) memory limit;
        memory request \(240\,Mi\) (HNU) or \(200\,Mi\) (LNU). \\
        Traffic & Constant-arrival-rate at \(8\) requests/s to the API only.
        Each request executes a deterministic CPU-bound loop; the experiment
        does not generate memory load. \\
        Actual Usage policy & CPU and memory thresholds \(0.80\);
        \textit{missingRequestPolicy=Allow}; namespace
        \textit{actual-usage-exp}; label
        \textit{experiment=actual-usage-evictor}. \\
        Timing & \(30\,s\) stabilization; two consecutive CPU samples at least
        \(10\,s\) apart above threshold; \(30\,s\) post-event observation and
        analysis window \([0,30)\,s\). \\
        Experimental runs and reset & One independent run for each of H0, H1, L0,
        and L1. Each runner cleans previous state and rebuilds and validates
        the layout before its run. \\
        \bottomrule
    \end{tabular}
\end{table*}

Every run uses a \textit{workload-api} Deployment configured with one replica
Pod and exposed through a dedicated NodePort Service. A \textit{k6}
constant-arrival-rate workload sends approximately eight requests per second
only to this Service. Each request performs fixed CPU work. The API container
declares a \(50\,m\) CPU request, a \(1000\,m\) CPU limit, and a \(700\,Mi\)
memory limit; its memory request is \(240\,Mi\) in HNU and \(200\,Mi\) in LNU.
Other experiment Pods receive no foreground traffic.

The system stabilizes under load for \(30\,s\). The Kubernetes Metrics API is
then polled every \(10\,s\), and the Descheduler is triggered only after two
consecutive samples place the API CPU usage above the configured
\(0.80\) request-relative threshold. The decision-time CPU ratios are
approximately \(10.02\) in H1 and \(10.10\) in L1.

Immediately before the single-shot Descheduler Job starts, the experiment
records an event timestamp. Load generation and Pod lifecycle observation
continue for a \(30\,s\) post-event period. Application metrics are calculated
from requests completed in the event-aligned interval \([0,30)\,s\).

Table~\ref{tab:experimental_setup} summarizes values recorded by the runners,
policies, plans, and run artifacts. H0, H1, L0, and L1 are each one independent
run, not repeated trials. Before each run, the runner deletes prior experiment
state and recreates and validates the intended layout.

\subsection{Comparison Systems and Layouts}

The paired layouts are designed to compare filter behavior while holding the
workload, host strategy, traffic, and intended initial placement constant.
Each pair differs in whether the Actual Usage Evictor is registered after the
Default Evictor. Table~\ref{tab:comparison_systems} summarizes the four
configurations.

\begin{table*}[t]
    \centering
    \caption{H0, H1, L0, and L1 Configurations}
    \label{tab:comparison_systems}
    \small
    \renewcommand{\arraystretch}{1.2}
    \begin{tabular}{@{}l>{\raggedright\arraybackslash}p{2.4cm}
        >{\raggedright\arraybackslash}p{1.5cm}
        >{\raggedright\arraybackslash}p{4.3cm}
        >{\raggedright\arraybackslash}p{2.3cm}
        >{\raggedright\arraybackslash}p{3.5cm}@{}}
        \toprule
        \textbf{Run} & \textbf{Host / scheduler} &
        \textbf{Filter} & \textbf{Intended initial layout} &
        \textbf{Thresholds} & \textbf{Outcome} \\
        \midrule
        H0 & HNU / MostAllocated & Disabled &
        Three dense destinations; three one-Pod sources, with API alone on one source. &
        HNU CPU/memory \(40\%\) & Evictions, active Nodes, API identity,
        failures, failure span. \\
        H1 & HNU / MostAllocated & Enabled &
        Same intended layout as H0. &
        HNU \(40\%\); Actual Usage CPU/memory \(0.80\) & Same as H0 plus
        Actual Usage blocks. \\
        L0 & LNU / LeastAllocated & Disabled &
        Three underutilized destinations; three overutilized sources; API and
        idle fallback share one source. &
        LNU low \(20\%\), target \(50\%\) & Evictions, request deviation, API
        identity, failures, failure span. \\
        L1 & LNU / LeastAllocated & Enabled &
        Same intended layout as L0. &
        LNU \(20/50\%\); Actual Usage CPU/memory \(0.80\) & Same as L0 plus
        Actual Usage blocks. \\
        \bottomrule
    \end{tabular}
\end{table*}

\textit{HighNodeUtilization} (HNU) uses \textit{MostAllocated} scoring to pack
replacement Pods. In H0, the strategy can evict all three source Pods and
empty their Nodes. In H1, protecting the API prevents worker-6 from becoming
empty, but the two idle source Pods remain eligible.

\textit{LowNodeUtilization} (LNU) uses the default
\textit{LeastAllocated} scheduler. In both runs, the strategy attempts one
movement from each overutilized source. On worker-6 the busy API is considered
before an idle fallback. Therefore, L1 can reject the API and continue with
the fallback without reducing the successful eviction count.

\subsection{Metrics}

The Actual Usage Block Count \(B\) counts candidate rejections caused by a CPU
or memory ratio reaching its threshold. HTTP Failure Rate \(F\) is

\[
F=
\frac{N_{\mathrm{fail,post}}}
     {N_{\mathrm{completed,post}}},
\]

where a failure is a completed non-\(2xx\) response in the common post-event
window. Observed Failure Span \(S_f\) is the elapsed time from the first to the
last failed response:

\[
S_f=
\begin{cases}
t_{\mathrm{last\ fail}}-t_{\mathrm{first\ fail}},
    & N_{\mathrm{fail,post}}>0,\\
0, & N_{\mathrm{fail,post}}=0.
\end{cases}
\]

This quantity describes the temporal spread of observed failures, not the full
period of service unavailability or time to recovery. It is zero both when no
failure is observed and when only one failure occurs.

The evaluation also records successful eviction count, evicted Pod identity,
and whether the original API instance remains. For HNU, placement progress is
the reduction in active Nodes hosting non-DaemonSet application workloads.
For LNU, it is the reduction in the population deviation of normalized CPU
and memory requests across Nodes. Successful-request tail latency is treated
as secondary because individual runs contain transient bursts before and
after the event; failure rate and observed failure span are the primary
availability evidence.

\section{Results and Discussion}
\label{sec:results}

\begin{table*}[t]
    \centering
    \caption{Consolidated Actual Usage Evictor Results}
    \label{tab:consolidated_results}
    \small
    \renewcommand{\arraystretch}{1.18}
    \begin{tabular}{@{}lccccp{6.0cm}@{}}
        \toprule
        \textbf{Run} & \textbf{Evictions} & \textbf{Blocks} &
        \textbf{Failure rate} & \textbf{Failure span} &
        \textbf{Strategy-specific outcome} \\
        \midrule
        H0 & 3 & 0 & \(40/240=16.67\%\) & \(4.836\,s\) &
        Active Nodes \(6\to3\); busy API evicted. \\
        H1 & 2 & 1 & \(0/240=0\%\) & \(0\,s\) &
        Active Nodes \(6\to4\); original API retained. \\
        L0 & 3 & 0 & \(40/232=17.24\%\) & \(4.873\,s\) &
        Memory deviation reduced \(85.7\%\); busy API evicted. \\
        L1 & 3 & 1 & \(0/240=0\%\) & \(0\,s\) &
        Same \(85.7\%\) reduction via idle fallback; original API retained. \\
        \bottomrule
    \end{tabular}
\end{table*}

\subsection{HighNodeUtilization}

H0 evicts all three source Pods, including the busy API, and reduces the
active-Node count from six to three. H1 allows both idle source Pods but
observes an API CPU ratio of approximately \(10.02\), blocks that candidate,
and reduces the active-Node count from six to four. The filter therefore does
not disable consolidation; it trades one additional reclaimable Node for
preserving the active service instance.

H0 completes \(240\) responses in the common window and returns \(40\) HTTP
\(503\) responses, a \(16.67\%\) failure rate. Those failures have an observed
span of \(4.836\,s\). H1 completes \(240\) responses without failure.

\subsection{LowNodeUtilization}

L0 evicts the busy API and one idle Pod from each of the other sources. L1
observes an API CPU ratio of approximately \(10.10\), rejects that candidate,
then permits the idle fallback on the same Node. Both configurations complete
three successful evictions and produce the same final request distribution.

Memory request deviation falls from \(0.21639\) to \(0.03091\) in both runs,
a reduction of \(0.18548\), or \(85.7\%\). CPU deviation is unchanged because
the layout creates imbalance through memory requests while candidate CPU
requests are equal. L0 completes \(232\) responses and returns \(40\) HTTP
\(503\) responses, a \(17.24\%\) failure rate with an observed failure span of
\(4.873\,s\). L1 completes \(240\) responses with no failure.

\subsection{Cross-Strategy Interpretation}

The treatment changes the decision for the busy API without making idle Pods
ineligible. The placement cost is determined by what the host strategy can do
next. HNU has no substitute on the protected one-Pod source, so that Node
remains active and consolidation is partial. LNU has an idle candidate on the
same source; it substitutes that candidate and preserves the complete measured
balancing result. Runtime protection and placement progress are therefore not
inherently opposed. They conflict only when the host strategy lacks a
compatible idle alternative.

\subsection{Limitations}

The experiments use controlled six-Node layouts, synthetic CPU-bound demand,
a single API replica, one traffic pattern, one threshold value, and one
independent run per configuration. The paired layouts support a direct
behavioral comparison but do not represent every workload mix or failure
recovery mechanism. A replicated service might mask instance replacement,
while stateful, memory-bound, or long-running requests could exhibit a
different failure profile.

The \(0.80\) threshold is evaluated as a decision setting, not optimized.
Metrics API observations are sampled rather than continuous and can lag rapid
load changes. Finally, only two upstream strategies are evaluated. They cover
opposing consolidation and redistribution objectives, but broader strategy
and scheduler combinations are needed before generalizing the placement-cost
results to production clusters.

\clearpage
